\begin{longtable}{@{}L{1.05cm} L{3.2cm} L{10.1cm}@{}}
\caption{Điểm yếu phương pháp luận nhóm phát hiện khi đọc toàn văn. Đây là cơ sở để nhóm khẳng định các con số công bố trong y văn không so sánh trực tiếp được với nhau.}\label{tab:diemyeu}\\
\toprule \textbf{Mã} & \textbf{Công trình} & \textbf{Hạn chế} \\
\midrule\endfirsthead
\multicolumn{3}{@{}l}{\footnotesize\itshape Bảng \thetable{} (tiếp theo)}\\
\toprule \textbf{Mã} & \textbf{Công trình} & \textbf{Hạn chế} \\
\midrule\endhead
\midrule\multicolumn{3}{r@{}}{\footnotesize\itshape tiếp trang sau}\\\endfoot
\bottomrule\endlastfoot
p01 & {\scriptsize Zhong 2018} & {\scriptsize Nghiêm trọng nhất, chỉ số F1 77,85\% không phải là F1 dò nhịp theo chuẩn CinC 2013 mà là F1 của bài toán phân loại cửa sổ 100 ms trên tập đã cân bằng lớp nhân tạo, trong khi toàn bộ văn liệu về sau (kể cả Mohebbian 2022 và Basak 2024 trong bảng so sánh của họ) đều trích con số này như F1 dò fQRS — một lỗi trích dẫn hệ thống của cả lĩnh vực. Không có hậu xử lý lấy đỉnh và không áp thời gian trơ nên một phức bộ QRS thật có thể sinh ra nhiều cửa sổ dương liên tiếp mà không bị tính là FP; chỉ có một lần chia cố định với 7 bản ghi kiểm thử, không lặp lại,…} \\[2pt]
p03 & {\scriptsize Mohebbian 2022} & {\scriptsize Nghiêm trọng nhất, với NI-FECG và CinC Challenge, tín hiệu fECG đích dùng để huấn luyện được tự sinh từ chuỗi R-R đáp án bằng wavelet Daubechies: mạng được dạy để tạo ra một tín hiệu mà các đỉnh nằm đúng vị trí đáp án, sau đó lại lấy đỉnh trên chính tín hiệu đó bằng Pan-Tompkins rồi so với đáp án, nên nếu có bất kỳ rò rỉ dữ liệu nào giữa fold huấn luyện và fold kiểm thử (điều bài không chứng minh là không có) thì con số 99,3\% sẽ hoàn toàn vô nghĩa. Bài không nêu cách chia 4 fold theo bản ghi hay theo cửa sổ ngẫu nhiên — đây là thông tin bắt buộc phải…} \\[2pt]
p04 & {\scriptsize Basak 2024} & {\scriptsize Nghiêm trọng nhất, bài chỉ nói "5-fold cross-validation" mà không nêu cách chia fold, trong khi có một dấu hiệu suy luận rất mạnh cho thấy fold được chia ngẫu nhiên theo đoạn chứ không theo chủ thể: ở Bảng 5, thống kê của chính dữ liệu thật (nhãn chuẩn) gần như y hệt nhau ở cả 5 fold, với muy\_RR = 473,19 / 473,18 / 473,19 / 473,18 / 473,20 ms và muy\_HR = 128,52 / 128,52 / 128,51 / 128,52 / 128,53 bpm, mà nếu mỗi fold kiểm thử chứa những thai nhi khác nhau thì nhịp tim thật trung bình phải khác nhau rõ rệt vì mỗi thai nhi có nhịp riêng — đây là rò rỉ dữ…} \\[2pt]
p05 & {\scriptsize Arash 2023} & {\scriptsize Bài không nêu số tham số, kích thước mô hình và số FLOPs nên không thể tái lập chính xác kiến trúc cũng như không so sánh được ngân sách tính toán, đồng thời tồn tại nhiều mâu thuẫn nội tại: mục IV-C ghi 6 bộ lọc độ dài "1, 3, 5, 7, 9 và 11" trong khi mục IV-E ghi "1, 3, 5, 7, 9 và 13"; dòng trung bình Bảng II ghi Se 86,23\% và PPV 95,75\% nhưng F1 tương ứng phải khoảng 90,72\% chứ không phải 93,43\%, và trung bình của 97,7\% với 95,3\% cũng không thể ra 93,43\%; Hình 9(a) ghi F1 tốt nhất 98,37\% tại 4 giây trong khi Bảng I ghi 99,03\% cho cùng cấu hình. Việc…} \\[2pt]
p06 & {\scriptsize Lin 2025} & {\scriptsize Rò rỉ dữ liệu tiềm ẩn nghiêm trọng trên ADFECGDB: mô hình được tinh chỉnh trên một phần ADFECGDB rồi báo cáo kết quả trên cả 5 bản ghi của chính bộ đó mà không ghi rõ cách chia, nên con số 99,17\% không được coi là kết quả ngoài mẫu — chính bài p07 (Asadi 2025) đã chỉ đích danh điều này khi viết rằng "một phần bộ ADFECGDB đã được dùng để huấn luyện trong các mô hình này, điều nhiều khả năng đã góp phần làm tăng hiệu năng của chúng trên bộ dữ liệu này", và Attention R2W-Net được nêu tên trong nhóm đó. Bài không nêu số tham số, kích thước mô hình hay số…} \\[2pt]
p07 & {\scriptsize Amirreza 2025} & {\scriptsize Kết quả tuyệt đối trên ADFECGDB (F1 96,52\%) thấp hơn đáng kể các bài có tinh chỉnh — chính nhóm tác giả thừa nhận thua TCGAN, CBLS-CycleGAN, Improved CycleGAN, DAW-Net, Attention R2W-Net và dual-branch Transformer-CNN trên bộ này, tuy lập luận (đúng) rằng đó là do khác biệt quy trình huấn luyện chứ không phải do mô hình kém; ngoài ra mô hình rất nặng, 9,22 triệu tham số và 111,79 MB cho một bài toán một kênh nên không thể triển khai trên thiết bị đeo, và bài báo cáo FLOPs nhưng không đo độ trễ thực tế (ms) trên bất kỳ CPU nào nên không biết có chạy…} \\[2pt]
p09 & {\scriptsize Xu 2026} & {\scriptsize Lỗ hổng tái lập nghiêm trọng nhất là bài không nêu dung sai khớp nhịp tính bằng ms, khiến mọi con số F1 không thể đối chiếu với bất kỳ công trình nào khác; bài cũng không viết công thức F1, cột "Accuracy" ở Bảng 8 ghi 0,0172 và 0,0152 cho DaISy là vô nghĩa và không định nghĩa được, còn Hình 11 ghi độ rộng QRS 28 ms với sóng T chỉ cách S 16 ms - theo suy luận của tôi là bất khả thi về sinh lý (QRS thai bình thường khoảng 50-60 ms) nên rất có thể đó là dao động của bộ lọc 10-60 Hz bị đọc nhầm thành sóng P và T. Cỡ mẫu đánh giá quá nhỏ: mỗi bản ghi DaISy…} \\[2pt]
p10 & {\scriptsize Esmaeili 2025} & {\scriptsize Mâu thuẫn số học nghiêm trọng (suy luận của tôi, đã kiểm chứng bằng tính toán): từ Sensitivity, Specificity và Precision họ báo cáo, giải ngược ra tỷ lệ lớp dương là 77,6\% ở phương án 1 và 77,3\% ở phương án 5, tính lại Accuracy từ tỷ lệ đó cho 0,9836 và 0,9675 - khớp chính xác với con số bài ghi (0,9836 và 0,9679) - nhưng với 10 ô mỗi giây và nhịp thai 110-160 bpm thì chỉ khoảng 20-27\% số ô có QRS thai, không thể là 77\%, nên các chỉ số này được tính với lớp dương là ô KHÔNG có QRS thai, tức "Sensitivity 97,91\%" trong tóm tắt thực chất là độ nhạy của…} \\[2pt]
p11 & {\scriptsize Orvas 2025} & {\scriptsize Mất đối xứng nghiêm trọng giữa mô phỏng và thực tế: F-score tụt từ 99,8 trên mô phỏng xuống 77,8 trên CinC 2013 thật, mất 22 điểm, mà bản thân bài không nhấn mạnh dù dữ liệu của chính họ cho thấy rõ; Bảng I cũng đáng ngờ khi F-score gần như bất biến 99,79-99,84 trên toàn dải SNR từ -25 dB đến 0 dB trong khi PCC rơi từ 97,80 xuống 64,65 và PRD tăng từ 21,89 lên 80,39 - suy luận của tôi là dữ liệu fecgsyn có nhịp R hoàn toàn xác định và tuần hoàn nên mạng có thể học "nhịp điệu" của bộ sinh thay vì học hình dạng QRS, đúng kiểu ngộ nhận mà nhóm đã tự phát…} \\[2pt]
p12 & {\scriptsize Wahbah 2024} & {\scriptsize Mô hình dùng đồng thời 12 kênh bụng nên không thể so với mô hình đơn kênh, và nhãn chuẩn là đầu ra thuật toán BSSR chứ không phải ECG da đầu, nghĩa là mô hình chỉ học bắt chước một thuật toán khác mà không ai biết BSSR chính xác đến đâu — đây là điểm yếu nghiêm trọng nhất về mặt khoa học. Bài không có dung sai khớp nhịp, dùng accuracy trên tập mất cân bằng làm metric chính (từ w1 = 3,17 suy ngược ra lớp đỉnh R chỉ chiếm khoảng 15,8\% số đoạn), và F1 = 0,96/0,97 chắc chắn không phải F1 của lớp đỉnh R vì với Se = 0,858 thì giải 2·Se·PPV/(Se+PPV) = 0,97…} \\[2pt]
p13 & {\scriptsize Behar 2014} & {\scriptsize Chương 7 báo cáo F1 trong mẫu vì mọi tham số đều quét trên set-a rồi F1 cuối cùng cũng tính trên set-a, còn tập giấu chỉ được chấm bằng điểm Challenge kiểu RMS; nhãn chuẩn của CinC 2013 lại không đồng nhất (một phần từ ECG da đầu, một phần gán tay, một phần suy từ kênh bụng mẹ, một phần từ tham số mô phỏng) và chính tác giả thừa nhận một số lỗi nhãn vẫn còn tồn tại, bên cạnh việc loại 7 bản ghi khỏi set-a bằng quan sát mắt thường sau khi đã nhìn dữ liệu. Không có bất kỳ con số chi phí tính toán nào trong 249 trang; chương chỉ số chất lượng tín hiệu…} \\[2pt]
p14 & {\scriptsize Behar 2014} & {\scriptsize Con số F1 = 96,0\% là trong mẫu vì mọi tham số quét trên set-a rồi F1 cũng báo cáo trên set-a, tập giấu chỉ có điểm Challenge kiểu RMS; thêm vào đó bài loại 7/75 bản ghi vì nhãn tham chiếu không chính xác xác định bằng quan sát mắt thường — quyết định chủ quan và làm sau khi đã nhìn dữ liệu — rồi còn cắt 2 giây đầu và cuối mỗi đoạn. Bốn ngưỡng SMI 29 bpm, BCM 0,4, SMI > 26 và nhịp hằng số 143 bpm đều đo bằng thực nghiệm trên set-a mà không có cơ sở lý thuyết, còn FUSE-CHALL được chính tác giả gọi là làm lệch về phía hệ thống chấm điểm của Challenge nên…} \\[2pt]
p15 & {\scriptsize Andreotti 2016} & {\scriptsize Toàn bộ đánh giá thực hiện trên dữ liệu mô phỏng và chính tác giả thừa nhận F1 cao một phần là do sự đơn giản của mô hình lưỡng cực trong fecgsyn, với khoảng FQT mô phỏng thấp hơn khoảng sinh lý thật (200-340 ms), tức dữ liệu tổng hợp chưa đúng về mặt hình thái. Việc chọn kênh kiểu oracle khiến mọi con số chỉ là cận trên lý thuyết chứ không phải hiệu năng của một hệ thống tự động; nhóm AM tuy được tuyên là đơn kênh nhưng trên lâm sàng vẫn phải dán thêm điện cực ngực, còn nhóm BSS cần 8 kênh đồng thời nên hoàn toàn không phải bài toán đơn kênh. Bài…} \\[2pt]
p16 & {\scriptsize Niknazar 2013} & {\scriptsize Bài không có một con số F1/Se/PPV nào nên không thể tái lập về mặt số học và không thể đặt vào bảng so sánh dò đỉnh, đồng thời bài không nêu dung sai khớp nhịp vì không đo khớp nhịp. Tồn tại một vòng lặp con gà - quả trứng: par-EKS cần vị trí đỉnh R thai làm đầu vào để khớp tham số mô hình, trong khi chính tác giả viết rằng phần cốt tử của par-EKS đề xuất là việc dò đỉnh R — nghĩa là đây không phải bộ dò đỉnh mà là bộ tách/tăng cường tín hiệu, bắt buộc phải có một bộ dò đỉnh khác đi trước; tuyên bố đơn kênh cũng bị phá vỡ ở ca song thai khi nhãn chuẩn…} \\[2pt]
p17 & {\scriptsize Castillo 2018} & {\scriptsize Điểm yếu lớn nhất là việc lọc bản ghi bằng bác sĩ: chỉ dùng 26/75 bản ghi Challenge Set A (khoảng 35\%) và loại 3/20 kênh ADFECGDB, mà chính các bản ghi bị loại mới là các bản ghi khó, nên F1 98,07\% không phải hiệu năng trên toàn bộ Set A mà là hiệu năng trên một tập con đã được làm sạch thủ công — bài có ghi rõ điều này nhưng phần tóm tắt và bảng so sánh lại không mang theo cảnh báo đó. Không có tách theo chủ thể hay tách theo bản ghi nên con số ADFECGDB là con số trong mẫu, và con số tổng bộ 94,11\% (vốn vẫn đã loại 3 tín hiệu) bị lu mờ trước con số…} \\[2pt]
p19 & {\scriptsize Eleonora 2020} & {\scriptsize Bài không nêu dung sai khớp nhịp bằng ms nên không thể tái lập con số độ chính xác, lại dùng công thức Acc có chứa TN mà không định nghĩa TN cho bài toán phát hiện đỉnh, khiến kết quả khó diễn giải và khó đối chiếu với F1/Se/PPV của các công trình khác. Siêu tham số lambda, N và lựa chọn bộ phát hiện jqrs đều được tinh chỉnh trên chính bộ dữ liệu thật được báo cáo, trong khi bộ dữ liệu đó lại riêng tư, chỉ 20 thai phụ, tuổi thai 21-27 tuần và mỗi bản ghi chỉ 10 giây, nên không tái lập độc lập được và không suy rộng cho theo dõi liên tục hay giai đoạn…} \\[2pt]
p20 & {\scriptsize Gari 2014} & {\scriptsize Nhãn tham chiếu giai đoạn cuối được tạo bằng cách hợp nhất nhãn gốc với nhãn của chính các thuật toán dự thi nên một phần nhãn là sản phẩm của thuật toán chứ không phải nhãn chuẩn độc lập, và tác giả thừa nhận vẫn còn lỗi; bài cũng không nêu dung sai khớp nhịp nên không thể tái lập. Bản ghi chỉ dài 1 phút nên không đánh giá được độ ổn định dài hạn, các tiêu chí chấm điểm theo (nhịp/phút)\textasciicircum{}2 và ms không đối chiếu được với F1/Se/PPV của đa số bài báo hiện đại khiến kết quả cuộc thi gần như không so sánh được với y văn sau này, còn sự kiện E3 kém tin cậy…} \\[2pt]
p21 & {\scriptsize Adam 2020} & {\scriptsize Số đối tượng rất nhỏ, chỉ 10 thai phụ ở B1 và 12 sản phụ ở B2, nên phương sai giữa các lần gấp khi tách theo bản ghi sẽ rất lớn (đây chính là lý do sai số chuẩn của nhóm là ±4,37 điểm), còn bản ghi B2 chỉ dài 5 phút và bản ghi B2\_Labour\_03 mất tới 27,5\% tín hiệu nên không đánh giá được độ ổn định dài hạn. Nhãn tham chiếu của B1 mang tính vòng tròn vì được suy ra từ chính tín hiệu bụng bằng thuật toán tự động cộng chuyên gia sửa tay, không có nhãn chuẩn độc lập và chính tác giả thừa nhận nhãn này chỉ "chính xác ở mức độ giới hạn", thêm 130 phức bộ…} \\[2pt]
p22 & {\scriptsize Jose 2013} & {\scriptsize Toàn bộ là dữ liệu tổng hợp với mỗi tín hiệu chỉ 50 mẫu và không có bất kỳ kiểm chứng nào trên tín hiệu sinh lý thật, trong khi chi tiết cài đặt đồng điều bền vững được hẹn sang một bài chưa công bố ([23]) nên không tái lập được nguyên vẹn. Điểm số phụ thuộc việc đoán đúng số chu kỳ L trong cửa sổ và thí nghiệm 1 thử ba giá trị L rồi lấy kết quả tốt nhất ngay trên tập đánh giá; ngoài ra bài thua Lomb-Scargle ở dạng cosin và có xu thế, tức với tín hiệu gần hình sin thì phương pháp cổ điển vẫn tốt hơn. Ràng buộc M ≥ 2N cùng cửa sổ khoảng một chu kỳ làm…} \\[2pt]
p23 & {\scriptsize Henry 2017} & {\scriptsize Thí nghiệm chính (Bảng 1) không phải phân loại có giám sát mà là phân cụm K-medoids không giám sát trên toàn bộ 150 mẫu, con số báo cáo lại là lần tốt nhất trong 1.000 khởi tạo nên không có tập kiểm thử độc lập và con số mang tính lạc quan; mục 6.4.1 cùng phần PI của mục 6.4.2 không nêu giao thức chia dữ liệu (số fold, tỷ lệ train/test) nên không tái lập được; và tham số b của hàm trọng số lấy từ độ dài bền vững lớn nhất của tất cả các lần thử kể cả dữ liệu đánh giá nên có rò rỉ dữ liệu nhẹ. Tính ổn định chỉ đúng với 1-Wasserstein, còn Remark 6 khẳng…} \\[2pt]
p24 & {\scriptsize Mathieu 2020} & {\scriptsize Bài báo cáo cột "Max" là kết quả 10-fold tốt nhất trong 10 lần lặp rồi đối chiếu với con số một lần của đối thủ nên đây là so sánh không đối xứng và thiên lệch lạc quan (con số trung thực là cột "Mean"), đồng thời Bảng 6 và Hình 8 chọn siêu tham số (kích thước lưới, φ, op, tham số t) dựa trên độ chính xác trên tập kiểm thử. Bài không nêu số tham số huấn luyện, không nêu kích thước mô hình và không nêu độ trễ suy luận một mẫu mà chỉ có thời gian chạy tổng của một lần thí nghiệm; đồng thời thua xa các phương pháp phi-tô-pô trên các bộ có nhãn đỉnh (NCI1…} \\[2pt]
p25 & {\scriptsize Anirudh 2020} & {\scriptsize Rò rỉ dữ liệu ở GENEactiv (chia 75/25 theo khung trên 152 người) cộng với số tuyệt đối rất thấp (F1 chỉ 46–58\% cho nhận dạng hoạt động) khiến mọi chênh lệch vài điểm trên nền 54\% đều kém thuyết phục. Lợi ích của PI là biên: CNN 1D tăng từ 53,56 lên 54,41, tức 0,85 điểm, trong khi độ lệch chuẩn là 0,21–0,31 và bài không nêu p-value cho phần chuỗi thời gian; chính tác giả viết rằng việc hợp nhất PI chỉ cải thiện "marginally". Bài không báo cáo số tham số của bất kỳ mạng nào nên không thể biết phần cải thiện đến từ thông tin tô-pô hay chỉ từ việc thêm…} \\[2pt]
p26 & {\scriptsize Renata 2023} & {\scriptsize Toàn bộ thí nghiệm chính là dữ liệu tổng hợp trong R\textasciicircum{}2 và R\textasciicircum{}3, dữ liệu thật duy nhất là ảnh lá cây FLAVIA; không có chuỗi thời gian, không có tín hiệu sinh lý, không có dữ liệu 1 chiều. Các đường cơ sở bị làm yếu một cách hệ thống: MLP chỉ huấn luyện 2 epoch cho mỗi cấu hình, ML và NN chỉ nhận ma trận khoảng cách của 100 điểm lấy mẫu con trong khi PH nhìn toàn bộ 1.000 hoặc 5.000 điểm, và PointNet không được tăng cường dữ liệu dù bài gốc của nó có dùng; chính tác giả thừa nhận họ không nhằm chứng minh PH vượt trội so với state-of-the-art. Kết quả tính…} \\[2pt]
p27 & {\scriptsize Quoc 2019} & {\scriptsize Bài không báo cáo chi phí tính toán dưới bất kỳ dạng nào — không có số giây, số tham số hay dung lượng bộ nhớ — dù phải tính 10 filtration Vietoris-Rips cho mỗi chuỗi, chắc chắn rất chậm nhưng không biết chậm bao nhiêu. Các đường cơ sở 1NN, EE và LS không được chạy lại mà lấy số từ tài liệu tham khảo, trong khi họ có hạ mẫu với hệ số 2, 2, 3 và 2 cho bốn bộ mà không giải thích lý do, nên so sánh không hoàn toàn đồng nhất; thêm vào đó tau của single-delay được chọn bằng giá trị lớn nhất trên tập kiểm tra. Không có kiểm định thống kê nên chênh lệch 90,0…} \\[2pt]
p28 & {\scriptsize Eugene 2023} & {\scriptsize Bài hoàn toàn không có dữ liệu ECG, không có bệnh nhân và không có chỉ số y sinh nên mọi con số không chuyển thẳng sang bài toán fQRS; việc chọn đỉnh phổ làm thủ công bằng mắt và chưa được tự động hóa. Bộ trung vị lỗi dự báo của tôm hùm trùng khít với bộ của Lorenz (0,106 | 0,122 | 0,112) khiến phần kết quả định lượng chính không kiểm chứng được, và phần khảo sát nhiễu mâu thuẫn nội bộ (phần chữ nói năm mức nhưng liệt kê bốn, còn Hình 14 hiện sáu mức 40/37/30/27/20/17 dB). Không có kiểm định thống kê, không lặp seed, không khoảng tin cậy nên khẳng định…} \\[2pt]
p29 & {\scriptsize Yonglian 2023} & {\scriptsize Rò rỉ dữ liệu nghiêm trọng ở hai tầng — chia ngẫu nhiên 10 mảnh không tách bệnh nhân hay bản ghi, cộng với việc lớp "unacceptable" phần lớn là bản sao có nhiễu của chính các tín hiệu "acceptable" trong cùng bộ dữ liệu — nên con số 98,55\% gần như chắc chắn lạc quan quá mức; bài toán còn bị làm dễ đi vì 88,9\% lớp âm là nhiễu tổng hợp cộng ở −10 dB, khiến mạng chỉ cần học có nhiễu bw/em/ma/Gauss hay không chứ không phải học chất lượng lâm sàng thật, với chỉ 988 mẫu xấu thật. Nhãn 12 chuyển đạo không nhất quán theo thừa nhận của chính tác giả, quy trình tự…} \\[2pt]
p30 & {\scriptsize Meryll 1906} & {\scriptsize Bài có mâu thuẫn nội bộ về số bệnh nhân giữ lại để kiểm tra (5 hay 15) và về tổng số bệnh nhân (240 trong phần chữ, 241 khi cộng Bảng 1), đồng thời không có kiểm định thống kê cho tuyên bố TDA cải thiện dù chỉ đưa ra 10 cặp số thô, không tính trung bình, không tính độ lệch chuẩn và TDA thua ở 4 trên 10 thí nghiệm phát hiện. Bài không tái lập được vì không nêu số tham số, kích thước mô hình, độ trễ suy luận, kiến trúc chi tiết từng kênh, số epoch tổng, batch size hay thư viện học sâu, và cũng không nêu độ dài cửa sổ theo giây do cắt theo số nhịp ("as…} \\[2pt]
p31 & {\scriptsize Yu-Min 1908} & {\scriptsize Bài không nêu dải lọc thông dải cho ECG — toàn bộ khâu tiền xử lý chỉ được dẫn chiếu sang bài [50] — nên khâu quan trọng nhất theo thực nghiệm của nhóm là không tái lập được; bài cũng không nêu số tham số mô hình, độ trễ suy luận hay kích thước mô hình, chỉ có thời gian huấn luyện 5-7 phút và thời gian chạy Ripser. Hiệu năng tuyệt đối rất thấp: F1 phát hiện trạng thái Thức chỉ 0,41-0,45, Kappa 0,24-0,32 (mức đồng thuận yếu), độ chính xác dương tính chỉ 34-38\% nghĩa là cứ 3 epoch báo "Thức" thì 2 sai, và bảng số thô cho thấy FP dao động 126-175 epoch…} \\[2pt]
p32 & {\scriptsize Andy 2025} & {\scriptsize Bài không có bài toán dự đoán nào — không train/test, không mô hình, không độ chính xác/F1/AUC — nên chỉ chứng minh các chỉ số tô-pô khác nhau giữa các nhóm tuổi chứ không chứng minh chúng dự đoán được tuổi hay bệnh lý của một cá thể mới, đó là khoảng cách rất lớn giữa "có ý nghĩa thống kê" và "dùng được trong lâm sàng". Nghiêm trọng hơn là mâu thuẫn nội tại: phần Tóm tắt và Thảo luận viết rằng sơ sinh và nhũ nhi có số lượng và cường độ đặc trưng bền vững cao hơn (N1 cao), trong khi mục Kết quả 4.2 viết ngược lại rằng sơ sinh có N1 và PE1 thấp hơn có ý…} \\[2pt]
p33 & {\scriptsize Alperen 2021} & {\scriptsize Con số đẹp nhất là kết quả gộp đa kênh chứ không phải đơn kênh: 94,46\% và 81,35\% trên WESAD là gộp toàn bộ 10 kênh ngực và cổ tay (ACC, RESP, ECG, EDA, EMG, TEMP, BVP), trong khi riêng ECG một kênh chỉ đạt 88,70\% (hai lớp) và 76,27\% (ba lớp), còn trên DriveDB ECG một kênh chỉ 65,74\% và 53,32\% tức gần như vô dụng — ai trích dẫn "98,07\%" mà không nói rõ đó là tín hiệu hô hấp (RESP) chứ không phải ECG là trích dẫn sai. Có một lỗi phương pháp xử lý tín hiệu thực sự: hạ mẫu 700 Hz xuống 100 Hz bằng cách lấy mỗi phần tử thứ 7 mà không lọc chống chồng phổ…} \\[2pt]
\end{longtable}